% ---------------------------------------------------------------- temporal detector
\subsection{A better detector does not produce a better system}
\label{sec:os-temporal}

Section~\ref{sec:os-depth} located roughly ten AUROC points of unused signal and pointed
at time of day. Acting on that raises a question of method before one of performance: we
identified the temporal axis by looking at target labels, so a detector built on it is
selected with target information through the analyst rather than through the code. Scoring
it as Tier-1 would be the error we document in others, arrived at by a different route.

The test is whether a rule that sees only source data would have chosen the same axis.
Hold out each known activity from the source prototypes in turn, score that activity's own
\emph{source} windows as though they were novel, and record which dimensions detect them;
averaging over the three known classes gives per-dimension weights that use no target data
at all. The rule selects the temporal axis in both directions. \texttt{t\_time\_sin} and
\texttt{t\_time\_cos} take the largest weights, and the four temporal dimensions average
$1.58$ and $1.96$ against $0.74$ and $0.57$ for the nine activity statistics. The axis is
therefore selectable without target labels on this data, which is what distinguishes a
usable rule from an observation we happened to make with the target labels in front of us.

The scope of that claim needs stating plainly, because it is the strongest interpretive
step in this paper. What we have shown is that one source-only rule selects the same axis
that target labels identify, in two transfer directions between two apartments of one
dataset, for one choice of withheld activities. Whether leave-one-class-out weighting
recovers the discriminative axis in general --- in other buildings, other sensor layouts,
other label splits, or where novelty is not separated by time of day --- is untested here
and does not follow from these two agreements. The rule could equally be selecting
temporal features because cyclical time is the one part of this feature vector not derived
from sensor activations, in which case it would transfer, or because toileting and
cleaning happen at distinctive hours in these two particular homes, in which case it would
not. We cannot distinguish those with two buildings, and a reader should treat the design
rule as a hypothesis with source-only support rather than as an established method.

% source: src/exp30_temporal_detector.py -> results/exp30_temporal_detector.csv
\begin{table*}[t]
\centering
\caption{Detectors that isolate or upweight the temporal axis. $\Delta$ is a paired
bootstrap over the target stream against the $13$-d mean baseline, $2000$ resamples.
Variants marked [ph] select the axis using target
knowledge: they are DIAGNOSTIC upper bounds in the sense of
Section~\ref{sec:os-setup}, not deployable results. [t1] variants use source data only. H and the rejection rate are at the fixed $q=5\%$
source-calibrated threshold.}
\label{tab:ostemporal}
\small
\begin{tabular}{lccc|ccc}
\toprule
 & \multicolumn{3}{c|}{$\text{S}_1\!\to\!\text{S}_2$} & \multicolumn{3}{c}{$\text{S}_2\!\to\!\text{S}_1$} \\
Detector & AUROC & $\Delta$ & H / rej & AUROC & $\Delta$ & H / rej \\
\midrule
$13$-d mean [t1]        & $0.6392$ & --- & $0.4957$ / $70\%$ & $0.5253$ & --- & $0.4279$ / $43\%$ \\
$13$-d max-dim [t1]     & $0.6308$ & $-0.009$ & $\mathbf{0.5235}$ / $57\%$ & $0.5745$ & $+0.049$ & $0.4410$ / $54\%$ \\
LOCO-weighted [t1]      & $0.6373$ & ns & $0.4069$ / $74\%$ & $0.5888$ & $+0.064$ & $\mathbf{0.4825}$ / $55\%$ \\
Temporal $4$-d [ph]     & $0.6389$ & ns & $0.4685$ / $23\%$ & $\mathbf{0.6616}$ & $+0.136$ & $0.4986$ / $25\%$ \\
Time-of-day $2$-d [ph]  & $0.5915$ & $-0.047$ & $0.5053$ / $26\%$ & $0.6266$ & $+0.101$ & $0.4207$ / $16\%$ \\
\bottomrule
\end{tabular}
\end{table*}

As a detector the temporal axis works, and it works only where there is a shift to detect.
Entering $\text{S}_1$, every variant separates from the baseline: the source-only
weighting gains $0.064$ AUROC and the post-hoc temporal distance $0.136$. Entering
$\text{S}_2$, where the cross-building penalty is $1.6$ points, nothing improves and two
variants are significantly worse. This is the same asymmetry the adaptation methods
followed in Section~\ref{sec:os-methods}, now in a component that does not adapt anything.

As a system it changes nothing, and that is the more useful result. At the fixed
source-calibrated threshold the source-only weighting rejects $74\%$ of the forward stream
and its H falls to $0.4069$, below the $0.5176$ of the configuration in
Section~\ref{sec:os-deploy}. The better-separating score is the worse-calibrated one. The
$13$-d max-dim variant shows the converse: it records the best H in both directions while
its AUROC is \emph{significantly worse} in one of them, so its apparent advantage is a
rejection-rate effect of the kind Section~\ref{sec:os-maha} dissects rather than better
detection. We decline to recommend it for that reason.

What this establishes is a constraint rather than a method. Detection quality and
deployed open-set performance come apart here: a score can gain $0.064$ AUROC and lose
$0.11$ H at the threshold a deployment would actually use. Until a threshold can be
calibrated so that it transfers, improving the detector does not improve the system, and
the limiting factor in Table~\ref{tab:osthresh} is not the one anybody is working on.
